% GENERATED FILE. Edit canonical lessons, then run npm run generate:books.
% canonical-source-hash: fnv1a64:0f1069a916767c3b
% canonical-lessons: TE-C416-pathyapustakam, TE-C416-vyasam, TE-C416-grantham, TE-C416-pathyamsam, TE-C416-samiksincu

\chapter{Books and Essays}
\label{ch:te-books-and-essays}
\hlchaptermodality{416}

\begin{chapteropening}
\textbf{By the end of this chapter:} \emph{I can talk about textbooks, essays, treatises and syllabuses.}

Complete the last lesson of chapter 416: I can talk about textbooks, essays, treatises and syllabuses.
\end{chapteropening}

\section[\texorpdfstring{\te{పాఠ్యపుస్తకం} (pāṭhyapustakaṁ)}{పాఠ్యపుస్తకం (pāṭhyapustakaṁ)}]{\te{పాఠ్యపుస్తకం} (pāṭhyapustakaṁ) — a textbook}

\label{lesson:TE-C416-pathyapustakam}



\begin{quote}
Before the new one: say the Telugu for to debate, then the Telugu for to instruct.
\end{quote}

\subsection*{You'll want to know: \te{పాఠ్యపుస్తకం}}
\textbf{\te{పాఠ్యపుస్తకం}} — \emph{pāṭhyapustakaṁ} — \textquotedblleft{}a textbook\textquotedblright{}.

\begin{quote}
\te{కొత్త} \te{పాఠ్యపుస్తకాలు} \te{ఇంకా} \te{రాలేదు}.
\end{quote}

\begin{quote}
\emph{kotta pāṭhyapustakālu iṅkā rālēdu.}
\end{quote}

\begin{quote}
The new textbooks have not come yet.
\end{quote}

Plural: \te{పాఠ్యపుస్తకాలు}, \emph{pāṭhyapustakālu}.

\te{పాఠ్య}, to be read and taught, + \te{పుస్తకం}, a book.

\begin{cousinweb}
\textbf{\te{పాఠ్యపుస్తకం}} (\emph{pāṭhyapustakaṁ}) is built from Sanskrit \emph{pāṭhya}, \textquotedblleft{}to be recited, to be taught\textquotedblright{}, and \emph{pustaka}, \textquotedblleft{}a book\textquotedblright{} (Monier-Williams), the word in \te{పుస్తకం}.
\end{cousinweb}

\begin{grammarlens}[title={What you've built}]
One more word for giving reasons and opinions.
\end{grammarlens}

\subsection*{Your turn}
\begin{itemize}
\raggedright
  \item \emph{Say it:} \emph{pāṭhyapustakaṁ}
  \item \emph{Say it:} \emph{pāṭhyapustakaṁ}, once more
  \item \emph{Say it:} \emph{kotta pāṭhyapustakālu iṅkā rālēdu.} — The new textbooks have not come yet.
  \item \emph{Recall:} say the Telugu for to debate, then the Telugu for to instruct, then say \emph{pāṭhyapustakaṁ} again
\end{itemize}

\begin{tcolorbox}[breakable,colback=teal!4,colframe=teal!35!black,title={Before you move on}]
What does \te{పాఠ్యపుస్తకం} mean? (\textquotedblleft{}A textbook\textquotedblright{}.) Say it once more.
\end{tcolorbox}
\section[\texorpdfstring{\te{వ్యాసం} (vyāsaṁ)}{వ్యాసం (vyāsaṁ)}]{\te{వ్యాసం} (vyāsaṁ) — an essay; an article}

\label{lesson:TE-C416-vyasam}



\begin{quote}
Before the new one: say the Telugu for to instruct, then the Telugu for a textbook.
\end{quote}

\subsection*{You'll want to know: \te{వ్యాసం}}
\textbf{\te{వ్యాసం}} — \emph{vyāsaṁ} — \textquotedblleft{}an essay; an article\textquotedblright{}.

\begin{quote}
\te{ఆమె} \te{పత్రికలో} \te{ఒక} \te{వ్యాసం} \te{రాసింది}.
\end{quote}

\begin{quote}
\emph{āme patrikalō oka vyāsaṁ rāsindi.}
\end{quote}

\begin{quote}
She wrote an article in the magazine.
\end{quote}

Plural: \te{వ్యాసాలు}, \emph{vyāsālu}.

\te{వ్యాసరచన}, essay writing, a school exercise.

\begin{cousinweb}
\textbf{\te{వ్యాసం}} (\emph{vyāsaṁ}) is borrowed from Sanskrit \emph{vyāsa}, \textquotedblleft{}a detailed account\textquotedblright{} (CDIAL 12214), and also \textquotedblleft{}extension, diffusion\textquotedblright{} (Monier-Williams).
\end{cousinweb}

\begin{grammarlens}[title={What you've built}]
One more word for giving reasons and opinions.
\end{grammarlens}

\subsection*{Your turn}
\begin{itemize}
\raggedright
  \item \emph{Say it:} \emph{vyāsaṁ}
  \item \emph{Say it:} \emph{vyāsaṁ}, once more
  \item \emph{Say it:} \emph{āme patrikalō oka vyāsaṁ rāsindi.} — She wrote an article in the magazine.
  \item \emph{Recall:} say the Telugu for to instruct, then the Telugu for a textbook, then say \emph{vyāsaṁ} again
\end{itemize}

\begin{tcolorbox}[breakable,colback=teal!4,colframe=teal!35!black,title={Before you move on}]
What does \te{వ్యాసం} mean? (\textquotedblleft{}An essay; an article\textquotedblright{}.) Say it once more.
\end{tcolorbox}
\section[\texorpdfstring{\te{గ్రంథం} (granthaṁ)}{గ్రంథం (granthaṁ)}]{\te{గ్రంథం} (granthaṁ) — a scholarly book; a treatise}

\label{lesson:TE-C416-grantham}



\begin{quote}
Before the new one: say the Telugu for a textbook, then the Telugu for an essay.
\end{quote}

\subsection*{You'll want to know: \te{గ్రంథం}}
\textbf{\te{గ్రంథం}} — \emph{granthaṁ} — \textquotedblleft{}a scholarly book; a treatise\textquotedblright{}.

\begin{quote}
\te{ఆ} \te{గ్రంథం} \te{వంద} \te{సంవత్సరాల} \te{కిందట} \te{రాశారు}.
\end{quote}

\begin{quote}
\emph{ā granthaṁ vanda saṁvatsarāla kindaṭa rāśāru.}
\end{quote}

\begin{quote}
That treatise was written a hundred years ago.
\end{quote}

Plural: \te{గ్రంథాలు}, \emph{granthālu}.

\te{గ్రంథాలయం}, a library, which you know, is built on it.

\begin{cousinweb}
\textbf{\te{గ్రంథం}} (\emph{granthaṁ}) is borrowed from Sanskrit \emph{grantha}, \textquotedblleft{}a knot\textquotedblright{}, and later \textquotedblleft{}a book\textquotedblright{} (CDIAL 4350), from \emph{granth}, \textquotedblleft{}to tie, to string together\textquotedblright{}.
\end{cousinweb}

\begin{grammarlens}[title={What you've built}]
One more word for giving reasons and opinions.
\end{grammarlens}

\subsection*{Your turn}
\begin{itemize}
\raggedright
  \item \emph{Say it:} \emph{granthaṁ}
  \item \emph{Say it:} \emph{granthaṁ}, once more
  \item \emph{Say it:} \emph{ā granthaṁ vanda saṁvatsarāla kindaṭa rāśāru.} — That treatise was written a hundred years ago.
  \item \emph{Recall:} say the Telugu for a textbook, then the Telugu for an essay, then say \emph{granthaṁ} again
\end{itemize}

\begin{tcolorbox}[breakable,colback=teal!4,colframe=teal!35!black,title={Before you move on}]
What does \te{గ్రంథం} mean? (\textquotedblleft{}A scholarly book; a treatise\textquotedblright{}.) Say it once more.
\end{tcolorbox}
\section[\texorpdfstring{\te{పాఠ్యాంశం} (pāṭhyāṁśaṁ)}{పాఠ్యాంశం (pāṭhyāṁśaṁ)}]{\te{పాఠ్యాంశం} (pāṭhyāṁśaṁ) — a syllabus topic}

\label{lesson:TE-C416-pathyamsam}



\begin{quote}
Before the new one: say the Telugu for an essay, then the Telugu for a scholarly book.
\end{quote}

\subsection*{You'll want to know: \te{పాఠ్యాంశం}}
\textbf{\te{పాఠ్యాంశం}} — \emph{pāṭhyāṁśaṁ} — \textquotedblleft{}a syllabus topic\textquotedblright{}.

\begin{quote}
\te{పరీక్షలో} \te{అన్ని} \te{పాఠ్యాంశాల} \te{మీద} \te{ప్రశ్నలు} \te{వస్తాయి}.
\end{quote}

\begin{quote}
\emph{parīkṣalō anni pāṭhyāṁśāla mīda praśnalu vastāyi.}
\end{quote}

\begin{quote}
In the exam there are questions on every topic of the syllabus.
\end{quote}

Plural: \te{పాఠ్యాంశాలు}, \emph{pāṭhyāṁśālu}.

\te{పాఠ్య}, to be taught, + \te{అంశం}, a point.

\begin{cousinweb}
\textbf{\te{పాఠ్యాంశం}} (\emph{pāṭhyāṁśaṁ}) is built from Sanskrit \emph{pāṭhya}, \textquotedblleft{}to be taught\textquotedblright{} (Monier-Williams), and \emph{aṁśa}, \textquotedblleft{}a part\textquotedblright{} (CDIAL 2), the word in \te{అంశం}.
\end{cousinweb}

\begin{grammarlens}[title={What you've built}]
One more word for giving reasons and opinions.
\end{grammarlens}

\subsection*{Your turn}
\begin{itemize}
\raggedright
  \item \emph{Say it:} \emph{pāṭhyāṁśaṁ}
  \item \emph{Say it:} \emph{pāṭhyāṁśaṁ}, once more
  \item \emph{Say it:} \emph{parīkṣalō anni pāṭhyāṁśāla mīda praśnalu vastāyi.} — In the exam there are questions on every topic of the syllabus.
  \item \emph{Recall:} say the Telugu for an essay, then the Telugu for a scholarly book, then say \emph{pāṭhyāṁśaṁ} again
\end{itemize}

\begin{tcolorbox}[breakable,colback=teal!4,colframe=teal!35!black,title={Before you move on}]
What does \te{పాఠ్యాంశం} mean? (\textquotedblleft{}A syllabus topic\textquotedblright{}.) Say it once more.
\end{tcolorbox}
\section[\texorpdfstring{\te{సమీక్షించు} (samīkṣiñcu)}{సమీక్షించు (samīkṣiñcu)}]{\te{సమీక్షించు} (samīkṣiñcu) — to review; to assess}

\label{lesson:TE-C416-samiksincu}



\begin{quote}
Before the new one: say the Telugu for a scholarly book, then the Telugu for a syllabus topic.
\end{quote}

\subsection*{You'll want to know: \te{సమీక్షించు}}
\textbf{\te{సమీక్షించు}} — \emph{samīkṣiñcu} — \textquotedblleft{}to review; to assess\textquotedblright{}.

\begin{quote}
\te{మంత్రి} \te{పథకాల} \te{పనిని} \te{సమీక్షించారు}.
\end{quote}

\begin{quote}
\emph{mantri pathakāla panini samīkṣiñcāru.}
\end{quote}

\begin{quote}
The minister reviewed the work on the schemes.
\end{quote}

Past (I …): \te{సమీక్షించాను}, \emph{samīkṣiñcānu}.

The noun is \te{సమీక్ష}, a review.

\begin{cousinweb}
\textbf{\te{సమీక్షించు}} (\emph{samīkṣiñcu}) is built on Sanskrit \emph{samīkṣā}, \textquotedblleft{}thorough or close inspection\textquotedblright{} (Monier-Williams), with the Telugu verb ending -\te{ఇంచు}.
\end{cousinweb}

\begin{grammarlens}[title={What you've built}]
One more word for giving reasons and opinions.
\end{grammarlens}

\subsection*{Your turn}
\begin{itemize}
\raggedright
  \item \emph{Say it:} \emph{samīkṣiñcu}
  \item \emph{Say it:} \emph{samīkṣiñcu}, once more
  \item \emph{Say it:} \emph{mantri pathakāla panini samīkṣiñcāru.} — The minister reviewed the work on the schemes.
  \item \emph{Recall:} say the Telugu for a scholarly book, then the Telugu for a syllabus topic, then say \emph{samīkṣiñcu} again
\end{itemize}

\begin{tcolorbox}[breakable,colback=teal!4,colframe=teal!35!black,title={Before you move on}]
What does \te{సమీక్షించు} mean? (\textquotedblleft{}To review; to assess\textquotedblright{}.) Say it once more.
\end{tcolorbox}
